% !TEX root = ../../main.tex
% C3.7 — Giao thức RTSP và giả lập luồng camera (RÚT GỌN 2026-09-29; bản cũ ở report/archive/)
% Nguồn: RFC 2326 (RTSP 1.0), RFC 7826 (RTSP 2.0, bỏ RECORD/ANNOUNCE), RFC 3550 (RTP), RFC 6184 (RTP payload H.264, đồng hồ 90 kHz),
% Wiegand 2003 (H.264), tài liệu FFmpeg (-re, -stream_loop, stream copy), MediaMTX.
% Khớp app: scripts/rtsp.ps1 (FFmpeg -re -stream_loop -1 -i camN.mp4 -c copy -f rtsp -rtsp_transport tcp); workers/sources/rtsp.py (PyAV, tcp,
% timeout kết nối/đọc 8 s, kết nối lại tối đa 3 lần, chờ 0,25 s nhân đôi đến 2 s, mốc thời gian theo đồng hồ worker).
\section{Giao thức RTSP và giả lập luồng camera}
\label{sec:3.7}

\subsection{Video nén H.264}
\label{sec:3.7.1}

Camera IP mã hóa hình ảnh thành video nén trước khi truyền, vì một khung hình 1920$\times$1080 chưa nén chiếm khoảng 6~MB, tức khoảng 370~MB mỗi giây ở 60 khung hình mỗi giây. H.264~\cite{wiegand2003h264} là chuẩn nén phổ biến của camera giám sát và cũng là định dạng video của đề tài. \emph{Khung I} chỉ dùng dự đoán trong chính khung hình nên giải mã được độc lập; \emph{khung P} và \emph{khung B} chỉ lưu dịch chuyển và phần sai khác so với khung tham chiếu, nên chỉ giải mã được khi đã có khung tham chiếu. Vì vậy bước giải mã phải thực hiện trên mọi khung hình dù hệ thống chỉ phân tích một phần (Mục~\ref{sec:3.4}), và một bộ giải mã bắt đầu nhận luồng giữa chừng phải đợi tới khung I kế tiếp.

\subsection{Giao thức RTSP và RTP}
\label{sec:3.7.2}

RTSP (Real Time Streaming Protocol)~\cite{rfc2326} là giao thức \emph{điều khiển}: phía nhận xác định luồng bằng một địa chỉ dạng \texttt{rtsp://<máy chủ>:<cổng>/<đường dẫn>} và trao đổi các lệnh \texttt{DESCRIBE}, \texttt{SETUP}, \texttt{PLAY}, \texttt{TEARDOWN} để lấy mô tả, thiết lập, bắt đầu và kết thúc phiên. Dữ liệu được mang bởi RTP (Real-time Transport Protocol)~\cite{rfc3550}: video nén được chia thành các gói có số thứ tự và dấu thời gian; với H.264, dấu thời gian dùng đồng hồ 90~kHz~\cite{rfc6184}. Gói RTP có thể gửi qua UDP, độ trễ thấp nhưng có thể mất gói, hoặc xen kẽ trên kết nối TCP của phiên RTSP, không mất gói nhưng có thể đến chậm khi nghẽn. RTSP 2.0~\cite{rfc7826} đã bỏ các lệnh dùng để đẩy luồng lên máy chủ, nên các công cụ giả lập của đề tài dùng phiên bản 1.0.

\subsection{Giả lập camera bằng máy chủ phát luồng}
\label{sec:3.7.3}

Để hệ thống được kiểm tra với đúng giao thức của camera IP, mỗi tệp video được phát lại thành một luồng RTSP đóng vai trò một camera. Máy chủ phát luồng MediaMTX~\cite{mediamtx2026} cho phép đẩy luồng lên và đọc luồng xuống qua RTSP, mỗi luồng ở một đường dẫn riêng như \texttt{/cam1}. Công cụ FFmpeg~\cite{ffmpeg2026docs} đọc tệp và đẩy lên MediaMTX với ba tùy chọn: \texttt{-re} để phát theo đúng tốc độ khung hình gốc, \texttt{-stream\_loop -1} để lặp vô hạn như camera chạy liên tục, và \texttt{-c copy} để sao chép nguyên dữ liệu nén mà không mã hóa lại, nên luồng mang đúng dữ liệu H.264 của tệp gốc và gần như không tốn tài nguyên. Luồng giả lập giống camera thật ở chỗ cùng giao thức, cùng định dạng và \emph{không chờ} phía nhận: nếu phía nhận xử lý chậm, dữ liệu mới vẫn tiếp tục được phát. Khác biệt là nội dung lặp theo chu kỳ và hầu như không có mất gói hay mất kết nối như mạng camera thực tế.

\subsection{Tiếp nhận luồng RTSP trong hệ thống}
\label{sec:3.7.4}

Tiến trình xử lý nền đọc luồng qua thư viện PyAV, vốn dùng các thư viện giải mã của FFmpeg, với dữ liệu RTP xen kẽ trên TCP để tránh mất gói làm hỏng ảnh. Việc kết nối và chờ dữ liệu đều giới hạn 8 giây; khi luồng gián đoạn, hệ thống thử kết nối lại tối đa ba lần với thời gian chờ tăng dần từ 0,25 giây đến 2 giây, sau đó ghi nhận lượt xử lý là lỗi. Mỗi khung hình được gắn mốc thời gian theo thời điểm hệ thống nhận được, thay vì dấu thời gian RTP của luồng (Mục~\ref{sec:3.3.5}). Vì tốc độ xử lý chậm hơn nhiều so với tốc độ phát, hệ thống không thể phân tích liên tục mọi camera; cách phân chia thời gian xử lý giữa các camera được trình bày ở Mục~\ref{sec:5.2.1}.
